% Part: normal-modal-logic
% Chapter: tableaux
% Section: soundness

\documentclass[../../../include/open-logic-section]{subfiles}

\begin{document}

\olfileid{nml}{tab}{sou}

\olsection{\Log{K} لئی صحت مندی}

\begin{editorial}
  صحت مندی دا ایہ ثبوت کلاسیکی قضیاتی منطق والا ثبوت دوبارہ ورتدا اے،
  یعنی ہر گل مڈھ توں ثابت کردا اے۔ جے صحت مندی دا اپنے آپ وچ پورا
  ثبوت چاہیدا اے تاں ایہ ٹھیک اے۔ جے تسی عام تابلو دی صحت مندی پہلاں
  ویکھ چکے او، تاں ایہ کجھ دہرایا ہویا لگے گا۔ اگے چل کے ایس اپنے آپ
  وچ پورے روپ تے غیر امکانی صورت اُتے ٹکے روپ وچوں چُن سکنا ممکن
  بناؤن دی وِونت اے۔
\end{editorial}

\begin{explain}
  اگیتے والے !!{tableau}s دی صحت مندی وکھاؤن لئی اسیں ثابت کرنا اے کہ جے
  \[
  \sFmla{\True}{!B_1}[1], \dots, \sFmla{\True}{!B_n}[1], \sFmla{\False}{!A}[1]
  \]
  دا اک بند !!{tableau} اے، تاں $!B_1, \dots, !B_n \Entails !A$۔
  اُلٹے نکار والا بیان ثابت کرنا سوکھا اے۔ جے کسے $\mModel{M}$ تے
  جہان~$w$ لئی، $\mSat{M}{!B_i}[w]$ سارے $i=1$, \dots,~$n$ لئی ہووے
  پر $\mSat{M}{!A}[w]$ ہووے، تاں کوئی !!{tableau} بند نہیں ہو سکدا۔
  اجیہا ردّی نمونہ وکھاؤندا اے کہ !!{tableau} دے شروع والے مفروضے پورے
  ہون جوگ نیں۔ ثبوت دا طریقہ ایہ اے کہ جدوں !!a{tableau} دی کسے شاخ
  اُتے سارے اگیتے والے !!{formula}s پورے ہون جوگ ہون، تاں قاعدہ لاؤن نال
  گھٹ توں گھٹ اک ودھی ہوئی شاخ وی پورا ہون جوگ ہوندی اے۔ بند شاخاں
  پوریاں ہون جوگ نہیں، ایس لئی اگیتے والے !!{formula}s دے کسے پورا ہون
  جوگ سیٹ دے ہر !!{tableau} وچ گھٹ توں گھٹ اک کھلی شاخ ہووے گی۔

امکانی صورت وچ ایہ طریقہ ورتن لئی سانوں «پورا ہون جوگ» دی تعریف
  امکانی نمونیاں تے اگیتیاں تک ودھاؤنی پئے گی۔ اوہ تعریف ملن پچھوں
  ثبوت سِدھا اے۔
\end{explain}

\begin{defn}
  منّ لو $P$ اگیتیاں دا کوئی سیٹ اے، یعنی $P \subseteq (\PosInt)^*
  \setminus \{\emptyseq\}$، تے $\mModel{M}$ اک نمونہ اے۔
  فنکشن~$f\colon P \to W$،~$P$ دی $\mModel{M}$ وچ \emph{تعبیر} اے جے
  جدوں وی $\sigma$ تے $\sigma.n$ دونیں~$P$ وچ ہون، تاں
  $Rf(\sigma)f(\sigma.n)$ ہووے۔

اگیتیاں دے سیٹ~$P$ دی تعبیر دے لحاظ نال اسیں ایہ تعریف کر سکدے آں:
  \begin{enumerate}
  \item $\mModel{M}$، $\sFmla{\True}{!A}[\sigma]$ نوں اُدوں تے صرف اُدوں
    پورا کردا اے جدوں $\mSat{M}{!A}[f(\sigma)]$۔
  \item $\mModel{M}$، $\sFmla{\False}{!A}[\sigma]$ نوں اُدوں تے صرف اُدوں
    پورا کردا اے جدوں $\mSat/{M}{!A}[f(\sigma)]$۔
  \end{enumerate}
\end{defn}

\begin{defn}
  منّ لو $\Gamma$ اگیتے والے !!{formula}s دا اک سیٹ اے، تے $P(\Gamma)$
  اوہدے وچ آئے اگیتیاں دا سیٹ اے۔ جے $f$،~$P(\Gamma)$ دی $\mModel{M}$
  وچ تعبیر اے، تاں اسیں آکھدے آں کہ $\mModel{M}$،$\Gamma$ نوں~$f$ دے لحاظ نال پورا کردا اے، $\mSat{M}{\Gamma}[f]$، جے $\mModel{M}$،
  ~$\Gamma$ دے ہر اگیتے والے !!{formula} نوں~$f$ دے لحاظ نال پورا کردا اے۔
  $\Gamma$ \emph{پورا ہون جوگ} اے اُدوں تے صرف اُدوں جدوں کوئی
  نمونہ~$\mModel{M}$ تے کوئی تعبیر~$f$، $P(\Gamma)$ لئی، ہووے جیہدے لئی
  $\mSat{M}{\Gamma}[f]$۔
\end{defn}

\begin{prop}
  جے $\Gamma$ وچ $\sFmla{\True}{!A}[\sigma]$ تے
  $\sFmla{\False}{!A}[\sigma]$ دونیں ہون، کسے !!{formula}~$!A$ تے
  اگیتے~$\sigma$ لئی، تاں $\Gamma$ پورا نہیں ہو سکدا۔
\end{prop}

\begin{proof}
  ایسا نمونہ~$\mModel{M}$ تے تعبیر~$f$، $P(\Gamma)$ لئی، ہو ای نہیں سکدے
  جیہناں لئی $\mSat{M}{!A}[f(\sigma)]$ تے
  $\mSat/{M}{!A}[f(\sigma)]$ دونیں ہون۔
\end{proof}

\begin{thm}[صحت مندی]
  \ollabel{thm:tableau-soundness}
  جے $\Gamma$ دا کوئی بند !!{tableau} اے، تاں $\Gamma$ پورا نہیں ہو سکدا۔
\end{thm}

\begin{proof}
!!a{tableau} دی شاخ پورا ہون جوگ اُدوں تے صرف اُدوں اے جدوں اوہدے اُتے
آئے !!{signed formula}s دا سیٹ پورا ہون جوگ اے۔ !!a{tableau} نوں پورا
ہون جوگ آکھاں گے جے اوہدے وچ گھٹ توں گھٹ اک پورا ہون جوگ شاخ ہووے۔

اسیں ایہ وکھاؤنے آں کہ کسے پورا ہون جوگ !!{tableau} نوں اخذ دے
کسے قاعدے نال ودھائیے تاں ہر واری پورا ہون جوگ !!{tableau} ای ملدا اے۔
ایس نال قضیہ ثابت ہو جائے گا۔ ہر بند !!{tableau} اوہناں قاعدیاں نوں
اوس !!{tableau} اُتے لا کے ملدا اے جیہدے وچ صرف~$\Gamma$ دے مفروضے نیں۔
سو جے $\Gamma$ پورا ہون جوگ ہووے، تاں اوہدا ہر !!{tableau} پورا ہون جوگ
ہووے گا۔ پر بند !!{tableau} پورا نہیں ہو سکدا، کیوں جے اوہدیاں ساریاں
شاخاں بند ہوندیاں نیں، تے بند شاخاں پوریاں نہیں ہو سکدیاں۔

منّ لو ساڈے کول پورا ہون جوگ !!{tableau} اے، یعنی اجیہا !!a{tableau}
جیہدے وچ گھٹ توں گھٹ اک پورا ہون جوگ شاخ اے۔ اخذ دا قاعدہ لاؤن نال
یا تاں کسے شاخ اُتے !!{signed formula}s ودھدے نیں، یا شاخ دو شاخاں وچ
ونڈی جاندی اے۔ جے !!{tableau} وچ کوئی پورا ہون جوگ شاخ اے جیہڑی ایس
قاعدے نال نہیں ودھی، تاں ودھے ہوئے !!{tableau} وچ وی اوہ پورا ہون جوگ
شاخ رہندی اے۔ ایس لئی ودھیا تابلو پورا ہون جوگ اے۔ سو اسیں صرف اوہ صورت
ویکھنی اے جیہدے وچ قاعدہ کسے پورا ہون جوگ شاخ اُتے لگدا اے۔

منّ لو $\Gamma$ اوس شاخ دے !!{signed formula}s دا سیٹ اے، تے
$\sFmla{S}{!A}[\sigma] \in \Gamma$ اوہ !!{signed formula} اے جیہدے اُتے
قاعدہ لگدا اے۔ جے قاعدے نال شاخ نہیں ونڈی جاندی، تاں اسیں وکھانا اے
کہ ودھی ہوئی شاخ، یعنی قاعدے دے نتیجیاں سنے $\Gamma$، ہُن وی پورا ہون
جوگ اے۔ جے قاعدے نال شاخ دو حصیاں وچ ونڈی جاندی اے، تاں اسیں وکھانا
اے کہ ملن والیاں دونہاں شاخاں وچوں گھٹ توں گھٹ اک پورا ہون جوگ اے۔
\tagfalse{prvDiamond}
پہلاں اوہ اخذ ویکھدے آں جیہناں دا صرف اک مقدمہ اے۔
\begin{enumerate}
\item شاخ، قاعدہ $\TRule{\True}{\lnot}$ نوں
  $\sFmla{\True}{\lnot !B}[\sigma] \in \Gamma$ اُتے لا کے ودھدی اے۔
  تاں ودھی ہوئی شاخ اُتے !!{signed formula}s
  $\Gamma \cup \{\sFmla{\False}{!B}[\sigma]\}$ نیں۔ منّ لو
  $\mSat{M}{\Gamma}[f]$۔ خاص طور اُتے $\mSat{M}{\lnot !B}[f(\sigma)]$۔
  ایس لئی $\mSat/{M}{!B}[f(\sigma)]$، یعنی $\mModel{M}$،
  $\sFmla{\False}{!B}[\sigma]$ نوں~$f$ دے لحاظ نال پورا کردا اے۔
\item شاخ، قاعدہ $\TRule{\False}{\lnot}$ نوں
  $\sFmla{\False}{\lnot !B}[\sigma] \in \Gamma$ اُتے لا کے ودھدی اے۔ مشق۔
\item شاخ، قاعدہ $\TRule{\True}{\land}$ نوں
  $\sFmla{\True}{!B \land !C}[\sigma] \in \Gamma$ اُتے لا کے ودھدی اے۔
  ایس نال شاخ اُتے دو نویں !!{signed formula}s
  $\sFmla{\True}{!B}[\sigma]$ تے $\sFmla{\True}{!C}[\sigma]$ ملدے نیں۔
  منّ لو $\mSat{M}{\Gamma}[f]$، خاص کر کے $\mSat{M}{!B \land
    !C}[f(\sigma)]$۔ تاں $\mSat{M}{!B}[f(\sigma)]$ تے
  $\mSat{M}{!C}[f(\sigma)]$۔ ایس دا مطلب اے کہ $\mModel{M}$،
  $\sFmla{\True}{!B}[\sigma]$ تے $\sFmla{\True}{!C}[\sigma]$
  دونہاں نوں~$f$ دے لحاظ نال پورا کردا اے۔
\item شاخ، قاعدہ $\TRule{\False}{\lor}$ نوں
  $\sFmla{\False}{!B \lor !C} \in \Gamma$ اُتے لا کے ودھدی اے۔ مشق۔
\item شاخ، قاعدہ $\TRule{\False}{\lif}$ نوں
  $\sFmla{\False}{!B \lif !C}[\sigma] \in \Gamma$ اُتے لا کے ودھدی اے۔
  شاخ اُتے دو نویں !!{signed formula}s ملدے نیں،
  $\sFmla{\True}{!B}[\sigma]$ تے $\sFmla{\False}{!C}[\sigma]$۔
  منّ لو $\mSat{M}{\Gamma}[f]$، خاص کر کے
  $\mSat/{M}{!B \lif !C}[f(\sigma)]$۔ تاں
  $\mSat{M}{!B}[f(\sigma)]$ تے $\mSat/{M}{!C}[f(\sigma)]$۔
  ایس دا مطلب اے کہ $\mModel{M}, f$،
  $\sFmla{\True}{!B}[\sigma]$ تے $\sFmla{\False}{!C}[\sigma]$ دونہاں نوں
  پورا کردے نیں۔

\iftag{prvBox}{%
\item شاخ، قاعدہ $\TRule{\True}{\Box}$ نوں
  $\sFmla{\True}{\Box !B}[\sigma] \in \Gamma$ اُتے لا کے ودھدی اے۔
  \iftag{probBox}{مشق۔}{ایس نال شاخ اُتے نواں !!{signed
      formula}~$\sFmla{\True}{!B}[\sigma.n]$ ملدا اے، کسے
    $\sigma.n \in P(\Gamma)$ لئی (کیوں جے $\sigma.n$ پہلاں توں ورتیا
    ہویا ہونا چاہیدا اے)۔ منّ لو $\mSat{M}{\Gamma}[f]$، خاص کر کے
    $\mSat{M}{\Box !B}[f(\sigma)]$۔ جدوں $f$ اگیتیاں دی تعبیر اے تے
    $\sigma$, $\sigma.n \in P(\Gamma)$ دونیں ہون، تاں اسیں جاندے آں کہ
    $Rf(\sigma)f(\sigma.n)$۔ ایس لئی $\mSat{M}{!B}[f(\sigma.n)]$،
    یعنی $\mModel{M}, f$، $\sFmla{\True}{!B}[\sigma.n]$ نوں پورا کردے نیں۔}

\item شاخ، قاعدہ $\TRule{\False}{\Box}$ نوں
  $\sFmla{\False}{\Box !B}[\sigma] \in \Gamma$ اُتے لا کے ودھدی اے۔
  \iftag{probBox}{مشق۔}{ایس نال نواں !!{signed
      formula}~$\sFmla{\False}{!A}[\sigma.n]$ ملدا اے، جتھے $\sigma.n$
    شاخ لئی نواں اگیتا اے، یعنی $\sigma.n \notin P(\Gamma)$۔
    $\Gamma$ پورا ہون جوگ اے، ایس لئی کوئی $\Struct{M}$ تے تعبیر~$f$، $P(\Gamma)$ لئی، موجود اے جیہدے لئی $\Sat{M}{\Gamma}[f]$، خاص کر کے
    $\mSat/{M}{\Box !B}[f(\sigma)]$۔ اسیں وکھانا اے کہ
    $\Gamma \cup \{\sFmla{\False}{!B}[\sigma.n]\}$ پورا ہون جوگ اے۔
    ایس لئی اسیں $P(\Gamma) \cup \{\sigma.n\}$ دی تعبیر ایویں
    متعین کردے آں:

$\mSat/{M}{\Box !B}[f(\sigma)]$ اے، ایس لئی کوئی $w \in W$ اے
  جیہدے لئی $Rf(\sigma)w$ تے $\mSat/{M}{!B}[w]$۔ منّ لو $f'$، $f$ ورگا
  ای اے، بس $f'(\sigma.n) = w$۔ $f'(\sigma) = f(\sigma)$ تے
  $Rf(\sigma)w$ ہون کر کے، ایس لئی $Rf'(\sigma)f'(\sigma.n)$، سو $f'$،
  $P(\Gamma) \cup \{\sigma.n\}$ دی تعبیر اے۔ صاف اے کہ
  $\mSat/{M}{!B}[f'(\sigma.n)]$۔ $f(\sigma') = f'(\sigma')$
  ہر اگیتے $\sigma' \in P(\Gamma)$ لئی اے، ایس لئی
  $\mSat{M}{\Gamma}[f']$۔ سو $\mModel{M}, f'$،
  $\Gamma \cup \{\sFmla{\False}{!B}[\sigma.n]\}$ نوں پورا کردے نیں۔}
  }{}

\iftag{prvDiamond}{%
\item شاخ، قاعدہ $\TRule{\True}{\Diamond}$ نوں
  $\sFmla{\True}{\Diamond !B}[\sigma] \in \Gamma$ اُتے لا کے ودھدی اے۔
  \iftag{probDiamond}{مشق۔}{ایس نال نواں !!{signed
      formula}~$\sFmla{\True}{!A}[\sigma.n]$ ملدا اے، جتھے $\sigma.n$
    شاخ لئی نواں اگیتا اے، یعنی $\sigma.n \notin P(\Gamma)$۔
    $\Gamma$ پورا ہون جوگ اے، ایس لئی کوئی $\Struct{M}$ تے تعبیر~$f$، $P(\Gamma)$ لئی، موجود اے جیہدے لئی $\Sat{M}{\Gamma}[f]$، خاص کر کے
    $\mSat{M}{\Diamond !B}[f(\sigma)]$۔ اسیں وکھانا اے کہ
    $\Gamma \cup \{\sFmla{\True}{!B}[\sigma.n]\}$ پورا ہون جوگ اے۔
    ایس لئی اسیں $P(\Gamma) \cup \{\sigma.n\}$ دی تعبیر ایویں
    متعین کردے آں:

$\mSat{M}{\Diamond !B}[f(\sigma)]$ اے، ایس لئی کوئی $w \in W$ اے
  جیہدے لئی $Rf(\sigma)w$ تے $\mSat{M}{!B}[w]$۔ منّ لو $f'$، $f$ ورگا
  ای اے، بس $f'(\sigma.n) = w$۔ $f'(\sigma) = f(\sigma)$ تے
  $Rf(\sigma)w$ ہون کر کے، ایس لئی $Rf'(\sigma)f'(\sigma.n)$، سو $f'$،
  $P(\Gamma) \cup \{\sigma.n\}$ دی تعبیر اے۔ صاف اے کہ
  $\mSat{M}{!B}[f'(\sigma.n)]$۔ $f(\sigma') = f'(\sigma')$
  ہر اگیتے $\sigma' \in P(\Gamma)$ لئی اے، ایس لئی
  $\mSat{M}{\Gamma}[f']$۔ سو $\mModel{M}, f'$،
  $\Gamma \cup \{\sFmla{\True}{!B}[\sigma.n]\}$ نوں پورا کردے نیں۔}
\item شاخ، قاعدہ $\TRule{\False}{\Diamond}$ نوں
  $\sFmla{\False}{\Diamond !B}[\sigma] \in \Gamma$ اُتے لا کے ودھدی اے۔
  \iftag{probDiamond}{مشق۔}{ایس نال شاخ اُتے نواں !!{signed
      formula}~$\sFmla{\False}{!B}[\sigma.n]$ ملدا اے، کسے
    $\sigma.n \in P(\Gamma)$ لئی (کیوں جے $\sigma.n$ پہلاں توں ورتیا
    ہویا ہونا چاہیدا اے)۔ منّ لو $\mSat{M}{\Gamma}[f]$، خاص کر کے
    $\mSat/{M}{\Diamond !B}[f(\sigma)]$۔ جدوں $f$ اگیتیاں دی تعبیر اے تے
    $\sigma$, $\sigma.n \in P(\Gamma)$ دونیں ہون، تاں اسیں جاندے آں کہ
    $Rf(\sigma)f(\sigma.n)$۔ ایس لئی $\mSat/{M}{!B}[f(\sigma.n)]$،
    یعنی $\mModel{M}, f$، $\sFmla{\False}{!B}[\sigma.n]$ نوں پورا کردے نیں۔}
  }{}
\end{enumerate}
ہُن اوہ اخذ ویکھدے آں جیہناں دے دو مقدمے نیں۔
\begin{enumerate}
\item شاخ، قاعدہ $\TRule{\False}{\land}$ نوں
  $\sFmla{\False}{!B \land !C}[\sigma] \in \Gamma$ اُتے لا کے ودھدی اے۔
  ایس نال دو شاخاں بندیاں نیں۔ کھبی شاخ $\sFmla{\False}{!B}[\sigma]$
  راہیں تے سجی شاخ $\sFmla{\False}{!C}[\sigma]$ راہیں اگے ودھدی اے۔
  منّ لو $\mSat{M}{\Gamma}[f]$، خاص کر کے
  $\mSat/{M}{!B \land !C}[f(\sigma)]$۔ تاں
  $\mSat/{M}{!B}[f(\sigma)]$ یا $\mSat/{M}{!C}[f(\sigma)]$۔
  پہلی صورت وچ $\mModel{M}, f$، $\sFmla{\False}{!B}[\sigma]$ نوں
  پورا کردے نیں، یعنی کھبی شاخ پورا ہون جوگ اے۔ دوجی صورت وچ
  $\mModel{M}, f$، $\sFmla{\False}{!C}[\sigma]$ نوں پورا کردے نیں،
  یعنی سجی شاخ پورا ہون جوگ اے۔
\item شاخ، قاعدہ $\TRule{\True}{\lor}$ نوں
    $\sFmla{\True}{!B \lor !C}[\sigma] \in \Gamma$ اُتے لا کے ودھدی اے۔ مشق۔
\item شاخ، قاعدہ $\TRule{\True}{\lif}$ نوں
    $\sFmla{\True}{!B \lif !C}[\sigma] \in \Gamma$ اُتے لا کے ودھدی اے۔ مشق۔
\end{enumerate}
\end{proof}

\begin{prob}
\olref[nml][tab][sou]{thm:tableau-soundness} دا ثبوت پورا کرو۔
\end{prob}

\begin{cor}
\ollabel{cor:entailment-soundness}
جے $\Gamma \Proves !A$، تاں $\Gamma \Entails !A$۔
\end{cor}

\begin{proof}
  جے $\Gamma \Proves !A$، تاں کجھ $!B_1$, \dots, $!B_n \in
  \Gamma$ لئی، $\Delta = \{\sFmla{\False}{!A}[1], \sFmla{\True}{!B_1}[1],
  \dots, \sFmla{\True}{!B_n}[1]\}$ دا بند !!{tableau} اے۔
  اسیں وکھانا اے کہ $\Gamma \Entails !A$۔ منّ لو ایہ نہیں۔ تاں کسے
  $\mModel{M}$ تے $w$ لئی $\mSat{M}{!B_i}[w]$، $i=1$, \dots,~$n$
  لئی اے، پر $\mSat/{M}{!A}[w]$۔ منّ لو $f(1) = w$۔ تاں $f$،
  ~$P(\Delta)$ دی~$\mModel{M}$ وچ تعبیر اے، تے $\mModel{M}$،
  ~$\Delta$ نوں~$f$ دے لحاظ نال پورا کردا اے۔ پر
  \olref{thm:tableau-soundness} مطابق، $\Delta$ دا بند !!{tableau} ہون
  کر کے اوہ پورا نہیں ہو سکدا۔ ایہ تضاد اے۔ سو آخرکار
  $\Gamma \Proves !A$ ہونا لازمی اے۔
\end{proof}

\begin{cor}
\ollabel{cor:weak-soundness}
جے $\Proves !A$، تاں $!A$ سارے نمونیاں وچ سچ اے۔
\end{cor}

\end{document}
